\scanpage{30}
\uline{Équisingularité}

\deleted{\unsure{On}} Le contexte des ``espaces'' est celui des espaces topologiques, ou des espaces linéaires \inserted{(ou algébriques, ou \unsure{analytiques})} par morceaux. \deleted{ou des espaces algébriques \unsure{par|p.ex.}} \unsure{Les mots|Le mot} ``espace'' et ``flèche'' désignent des \deleted{\unsure{tels}} objets et morphismes du contexte choisi

Soient $X$ un espace, $Y$ un sous-espace loc. fermé. $X$ est dit \uline{équisingulier le long de $Y$} si $\forall y\in Y$, $\exists$ ouvert $X'$ de $X$, et espace pointé $(Z,a)$, tels que l'on ait un diagramme commutatif
\diagram{\begin{tikzcd}
X'\cap Y=Y'_{\unsure{a|u}} \arrow[r,hook] \arrow[d,"\unsure{\wr|2}"] & X' \arrow[d,"\unsure{\wr|2}"] \\
Y'\times\{a\} \arrow[r,hook] & Y'\times Z
\end{tikzcd}}

Plus généralement, soit $\mathcal X=(X_i)_{i\in I}$, un diagramme d'espaces ($I$ étant un type de diagramme, p.ex. défini par un ens. ordonné --- c'est le seul cas que nous aurons à envisager), soit $i_0\in I$ et $Y$ un sous-espace loc. fermé de $X_{i_0}$. On dit que $\mathcal X$ est \uline{équisingulier le long de $Y$} si, désignant par $J$ le sous-diagramme de $I$ formé de tous les $j\in I$ pour lesquels \inserted{ou bien $j=i_0$ ou bien} $\exists$ flèche $j\longrightarrow i_0$, $\forall y\in Y$ \unsure{peut|pourra} trouver un voisinage ouvert $X'_{i_0}$ de $y$ dans $X_{i_0}$, un diagramme \inserted{d'espaces} $(Z_j)_{j\in J}$, avec $Z_{i_0}$ muni d'un point $a\in Z_{i_0}$, et un \unsure{homomorphisme|homéomorphisme} des diagrammes (où $Y'=Y\cap X'_{i_0}$)
\diagram{\begin{tikzcd}
(Z_j\times Y')_{j\in J} \arrow[r] & (X_j)_{j\in J}
\end{tikzcd}}
tel que pour tout $j\in J$, $Z'_j\times Y' - X_j$ soit un plongement ouvert, et que $\exists$ diagramme comm.
\diagram{\begin{tikzcd}
Y' \arrow[r,hook] \arrow[d,no head,"\unsure{\wr|2}"] & X'_{i_0} \arrow[d,no head,"\unsure{\wr|2}"] \\
\{a\}\times Y' \arrow[r,hook] & Z_{i_0}\times Y'
\end{tikzcd}}
\mspage{31}
